\begin{textsample}{POS Dim 3 – ac – Score 53.00 – ac/ac\_inf\_18.txt}  \label{ex:f3_pos_022}
18. \textbf{CRIANÇAS} BEM \textbf{PEQUENAS} ( 1 ANO E 7 \textbf{MESES} A 3 ANOS E 11 \textbf{MESES} ) 18.1.
Campo de experiências “ Corpo, \textbf{gestos} e movimentos ” Com o corpo ( por meio dos sentidos, \textbf{gestos}, movimentos impulsivos ou intencionais, coordenados ou espontâneos ), as \textbf{crianças}, desde cedo, exploram o mundo, o espaço e os objetos do seu entorno, estabelecem relações, expressam -se, \textbf{brincam} e produzem conhecimentos sobre si, sobre o outro, sobre o universo social e cultural, tornando -se, progressivamente, conscientes dessa corporeidade.
Por meio das diferentes linguagens, como a música, a dança, o teatro, as \textbf{brincadeiras} de faz de conta, elas se comunicam e se expressam no entrelaçamento entre corpo, emoção e linguagem.
As \textbf{crianças} conhecem e reconhecem as sensações e funções de seu corpo e, com seus \textbf{gestos} e movimentos, identificam suas potencialidades e seus limites, desenvolvendo, ao mesmo tempo, a consciência sobre o que é seguro e o que pode ser um risco à sua integridade física.
Na Educação \textbf{Infantil}, o corpo das \textbf{crianças} ganha \textbf{centralidade}, pois ele é o \textbf{partícipe} privilegiado das práticas pedagógicas de \textbf{cuidado} físico, orientadas para a emancipação e a liberdade, e não para a submissão.
Assim, a \textbf{instituição} escolar precisa promover oportunidades ricas para que as \textbf{crianças} possam, sempre animadas pelo espírito \textbf{lúdico} e na interação com seus \textbf{pares}, explorar e vivenciar um amplo repertório de movimentos, \textbf{gestos}, olhares, \textbf{sons} e \textbf{mímicas} com o corpo, para descobrir variados modos de ocupação e uso do espaço com o corpo ( tais como sentar com apoio, rastejar, engatinhar, escorregar, caminhar \textbf{apoiando} -se em berços, mesas e cordas, saltar, escalar, equilibrar -se, correr, dar cambalhotas, alongar -se etc. ). 18.2.
Direitos de aprendizagem e desenvolvimento - Conviver com \textbf{crianças} e \textbf{adultos}, experimentando as possibilidades corporais para interagir e desenvolver o interesse sobre os \textbf{cuidados} com o corpo. - \textbf{Brincar} utilizando criativamente o repertório da cultura corporal e do movimento, expressando emoções, alegrias e ( re ) significando o mundo a sua \textbf{volta}. - Explorar os movimentos corporais através de \textbf{brincadeiras} e interações para expressar suas emoções e desenvolver sua autonomia. - Participar de atividades que envolvem práticas corporais como dança, música, teatro, artes circenses, manifestando encanto e satisfação, desenvolvendo a autonomia. - Expressar corporalmente emoções e representações tanto nas relações cotidianas como nas \textbf{brincadeiras}, artes visuais, dramatizações, danças, músicas, \textbf{escuta} e contação de histórias. - Conhecer -se por meio de diversas experiências de interações, \textbf{brincadeiras} e explorações com seu corpo, desenvolvendo a autonomia e a autoestima. 18.2.1.
Objetivo de aprendizagem e desenvolvimento : ( EI02CG01 ) - Apropriar -se de \textbf{gestos} e movimentos de sua cultura no \textbf{cuidado} de si e nos jogos e \textbf{brincadeiras}.
Aprendizagens específicas - Desenvolver, gradativamente, uma \textbf{percepção} integrada do próprio corpo. - Expressar -se de diferentes modos através dos movimentos do corpo. - Desenvolver, gradativamente, o controle sobre o próprio corpo em danças, jogos, \textbf{brincadeiras} e outros movimentos. - Desenvolver a capacidade de criar \textbf{gestos} e movimentos, a partir da observação de diferentes manifestações corporais do seu grupo. - Expressar suas vivências e experiências de \textbf{cuidados} de si e \textbf{interessar} -se em \textbf{escutar} os relatos dos \textbf{colegas}. - Reconhecer as possibilidades do seu corpo, expandindo seus movimentos nos diferentes espaços. - \textbf{Interessar} -se por compartilhar conhecimentos sobre jogos e \textbf{brincadeiras} e apreciar os dos seus \textbf{colegas}, ampliando seu repertório cultural. \textbf{Sugestões} de experiências - \textbf{Brincar} de faz de conta, utilizando como referência enredos, cenários e \textbf{personagens} do seu convívio social. - \textbf{Brincar} com jogos de sua cultura familiar e com os jogos de seus \textbf{colegas}. - Relatar práticas de \textbf{cuidados} de si em casa e \textbf{escutar} com atenção o relato dos \textbf{colegas}. - \textbf{Brincar} de \textbf{imitar} \textbf{gestos} e movimentos aprendidos com \textbf{colegas} e \textbf{adultos}. - \textbf{Brincar} de organizar os diversos espaços e materiais. - \textbf{Brincar} de \textbf{imitar} \textbf{gestos} e movimentos típicos dos profissionais da \textbf{instituição} e de sua comunidade local. - Cantar \textbf{imitando} os \textbf{gestos} ou seguir ritmos diferentes de músicas com movimentos corporais. - \textbf{Escutar} diferentes estilos de música ( local, regional e mundial ). - Participar de \textbf{brincadeiras} cantadas. - Dançar livremente. - \textbf{Imitar} movimentos e \textbf{gestos} a partir de apresentações artísticas assistidas. - \textbf{Imitar} movimentos na dança a partir do contato com diferentes gêneros musicais. - \textbf{Brincar} com argila. - \textbf{Brincar} de plantar, semear, \textbf{colher}, pescar, caçar. - \textbf{Brincar} deslocando -se e movimentando -se em espaços internos e externos da \textbf{instituição}. - \textbf{Brincar} movimentando os braços, as pernas, a cabeça, com orientação ou não do professor. - Comunicar -se através de \textbf{gestos}, expressões faciais, ritmos e movimentos corporais. - \textbf{Brincar} de \textbf{mímica}, dança, circo e outras possibilidades. - Cantar e dançar, em sintonia com o ritmo da música. - \textbf{Brincar} com diferentes materiais, como : bambolês, cordas, bolas, \textbf{pés} de lata e outros materiais. 18.2.2.
Objetivo de aprendizagem e desenvolvimento : ( EI02CG02 ) - Deslocar seu corpo no espaço, orientando -se por noções como em frente, atrás, no alto, embaixo, dentro, fora etc., ao se envolver em \textbf{brincadeiras} e atividades de diferentes naturezas.
Aprendizagens específicas - Reconhecer o espaço e movimentar -se, explorando -o de diferentes formas. - Desenvolver o interesse por fazer novas descobertas em seus deslocamentos e explorações nos diferentes espaços. - Desenvolver a orientação espacial nos deslocamentos, utilizando os movimentos corporais. - Desenvolver atitude de confiança e segurança nas próprias capacidades motoras, ao movimentar -se no espaço. - Desenvolver a autonomia ao explorar o espaço de diferentes formas e por diferentes perspectivas. - Desafiar -se a deslocar -se no espaço a partir das observações que faz, dos \textbf{adultos} e de seus \textbf{pares}. - Construir estratégias para a superação de desafios em seu deslocamento no espaço. - Apropriar -se do espaço a sua \textbf{volta}, reconhecendo onde se encontram seus pertences pessoais. \textbf{Sugestões} de experiências - Localizar um \textbf{brinquedo} e buscá -lo no espaço em que se encontra. - Localizar seus pertences no espaço para fazer uso conforme suas necessidades. - Saltar, correr, se arrastar explorando o espaço em que se encontra. - \textbf{Brincar} com os \textbf{colegas} de \textbf{esconder} e achar \textbf{brinquedos} e objetos no espaço. - Andar pelo espaço segurando objetos na \textbf{mão}. - \textbf{Imitar} seus \textbf{colegas} nas diferentes formas de exploração do espaço. - Participar de circuitos elaborados com diferentes materiais. - \textbf{Brincar} com objetos e \textbf{brinquedos} de materiais diversos, experimentando suas possibilidades : encher, apertar, morder, produzir \textbf{sons}, esvaziar, transvazar, empilhar, colocar dentro/fora, fazer afundar, flutuar, soprar, \textbf{montar}, construir, dentre outras. - \textbf{Brincar} em diferentes posições : deitado, em cima, em baixo, de lado, atrás, à frente e outras possibilidades. 18.2.3.
Objetivo de aprendizagem e desenvolvimento : ( EI02CG03 ) - Explorar formas de deslocamento no espaço ( pular, saltar, dançar ), combinando movimentos e seguindo orientações.
Aprendizagens específicas - Conhecer e reconhecer os diferentes espaços físicos da \textbf{instituição}, a fim de situar -se e deslocar -se de diferentes formas com segurança. - Reconhecer os diferentes espaços em que convive, a fim de compreender a funcionalidade de cada ambiente. - Adquirir, gradativamente, \textbf{equilíbrio}, ritmo e resistência nos movimentos, nas \textbf{brincadeiras}, danças, coreografias, dentre outros. - Ampliar o conhecimento sobre seu corpo no espaço, \textbf{demonstrando} domínio sobre ele, desafiando -se nos deslocamentos. - Ampliar \textbf{gestos} e movimentos a partir de ações já conhecidas, criando e reinventando outras cada vez mais complexas. - Desenvolver a corporeidade por meio de \textbf{brincadeiras}, músicas e danças, a partir de orientações verbais, visuais e de movimentos simples. - Desenvolver atitude de apreciação do movimento dos \textbf{adultos} e de seus \textbf{pares}, narrando, descrevendo e \textbf{demonstrando} nas danças, músicas e \textbf{brincadeiras}. - Utilizar as possíveis variações do movimento, como a velocidade, força, flexibilidade para se deslocar, \textbf{brincar} e jogar. \textbf{Sugestões} de experiências - Pegar o \textbf{brinquedo} quando solicitado pelo outro. - Correr para ver quem chega primeiro a um lugar marcado. - Correr, saltar ou rastejar -se, explorando espaços familiares e com obstáculos. - \textbf{Brincar} de esconde-esconde, pular dentro e fora de pneus, escorregar, balançar, dançar, rodopiar, andar, engatinhar, arrastar -se, pular, subir, descer, pendurar -se, saltar, equilibrar -se, perseguir, procurar e pegar. - Movimentar -se ao ritmo de músicas diversas. - Movimentar -se ao comando de seus \textbf{pares} e \textbf{adultos} ( ao \textbf{som} de palmas, estalar de \textbf{dedos}, \textbf{batidas} de \textbf{pés}, etc. ). 18.2.4.
Objetivo de aprendizagem e desenvolvimento : ( EI02CG04 ) - \textbf{Demonstrar} progressiva independência no \textbf{cuidado} do seu corpo.
Aprendizagens específicas - Reconhecer o próprio corpo e o dos \textbf{colegas}, demostrando cada vez mais \textbf{cuidado} consigo e com o outro. - Apropriar -se dos nomes das partes do corpo de forma contextualizada, em situações reais ou de faz de conta. - Adotar \textbf{hábitos} e \textbf{cuidados} saudáveis de \textbf{higiene} corporal. - Desenvolver, gradativamente, autonomia nos movimentos ao \textbf{cuidar} do corpo. - Expressar atitudes de \textbf{cuidado} e respeito consigo mesmo e com o outro. - Desenvolver, progressivamente, autonomia em situações de alimentação, \textbf{higiene} corporal, troca de roupa e em outras situações. - Desenvolver atitude de confiança e segurança ao solicitar \textbf{ajuda} do \textbf{adulto}, quando necessário. - Desenvolver o prazer da descoberta, experimentando novos alimentos. - Superar desafios frente a novas experiências com relação aos \textbf{cuidados} do seu corpo. - Desenvolver \textbf{hábitos} de \textbf{cuidado} de si, que envolvam ações de vestir -se, alimentar -se e outras. \textbf{Sugestões} de experiências - Alimentar -se, usar o vaso sanitário, limpar o nariz, lavar as \textbf{mãos}, escovar os dentes, tomar banho, ensaboando o corpo e os cabelos, pentear -se, autonomamente. - Calçar meias e sapatos, vestir -se com ou sem \textbf{ajuda}. - Solicitar \textbf{ajuda} para limpar -se ao usar o banheiro, quando necessário. - Alimentar -se, solicitando \textbf{ajuda} quando necessário. 18.2.5.
Objetivo de aprendizagem e desenvolvimento : ( EI02CG05 ) - Desenvolver progressivamente as habilidades manuais, adquirindo controle para desenhar, pintar, rasgar, folhear, entre outros.
Aprendizagens específicas - Desenvolver a coordenação motora ampla e fina para que tenha maior domínio no seu deslocamento e na atuação sobre os objetos. - Utilizar, progressivamente, as habilidades motoras para \textbf{manusear} diferentes objetos e materiais. - Desenvolver autocontrole de movimentos ao \textbf{manipular} materiais de diferentes tipos, tamanhos, espessuras, \textbf{texturas} e \textbf{pesos}. - Aprimorar a coordenação visio-motora fina, utilizando os movimentos de preensão para fazer suas \textbf{marcas} gráficas. - Desenvolver a \textbf{curiosidade} por conhecer novos objetos, seus usos e funções, \textbf{manuseando} -os em diferentes situações. - Desenvolver a \textbf{autoconfiança}, perseverança e interesse, frente aos desafios encontrados na manipulação de objetos e materiais. \textbf{Sugestões} de experiências - \textbf{Montar} \textbf{brinquedos}, pegar objetos e manuseá -los. - Carregar objetos controlando e equilibrando -os. - Construir \textbf{brinquedos} diversos com sucatas, com materiais da natureza e outros. - \textbf{Brincar} de cantar, dançar, desenhar, escrever, jogar \textbf{futebol}, jogar bola ao cesto, boliche, esconde-esconde, mapa do tesouro, estátua ou malabarista de circo. - Fazer \textbf{marcas} gráficas com giz de cera, canetas, \textbf{lápis}, pincéis grosso e fino, pincéis de rolinho, giz pastel e outros materiais. - Folhear páginas de \textbf{livros}, cadernos, revistas e outros. - \textbf{Brincar} com blocos de encaixe para construir objetos de diferentes tamanhos e formatos. - Usar, gradativamente, tesoura simples ( sem ponta ) para recortar. - Explorar as diferentes características de vários materiais, as possibilidades de manuseá -los e \textbf{brincar} com eles. 18.3.
Avaliação - \textbf{acompanhamento} do processo de aprendizagem e desenvolvimento das \textbf{crianças} bem \textbf{pequenas} : - Observação \textbf{criteriosa}, \textbf{registro} e análise quanto : - Às iniciativas em criar \textbf{gestos} e movimentos, a partir das manifestações corporais do seu grupo cultural. - Ao desempenho corporal nas explorações do espaço de diferentes formas e diferentes perspectivas. - À coordenação simultânea de movimentos e \textbf{gestos} de forma segura e equilibrada. - Às atitudes de \textbf{cuidados} pessoais com independência e autonomia. - Ao domínio e controle das habilidades manuais.

% matched lemmas: acompanhamento, adulto, ajuda, apoiar, autoconfiança, bater, brincadeira, brincar, brinquedo, centralidade, colega, colher, criança, criterioso, cuidado, cuidar, curiosidade, dedo, demonstrar, equilíbrio, esconder, escuta, escutar, futebol, gesto, higiene, hábito, imitar, infantil, instituição, interessar, livro, lápis, lúdico, manipular, manusear, marca, montar, mão, mês, mímico, par, partícipe, pequeno, percepção, personagem, peso, pé, registro, som, sugestão, textura, volta
\end{textsample}
